\renewcommand{\deutschertitel}{Methoden zur Positionsbestimmung}
\renewcommand{\englischertitel}{Position Estimation Methods}
\chapter[\texorpdfstring{\protect{\vspace{2pt}\englischertitel}}{\englischertitel}]{}
\kapitel{\deutschertitel}
\thispagestyle{empty}

\label{PositionEstimationMethods}

\begin{paracol}{2}[]

{\raggedright\huge\bfseries\sffamily \englischertitel \par\ } \\[1.8ex]

\switchcolumn

{\raggedright\huge\bfseries\sffamily \deutschertitel \par\ } \\[1.8ex]

\coleng

For a UAV to operate autonomously, it must maintain a state estimate—a real-time calculation of its position, velocity, and orientation. To estimate the position and movement of the vehicle, the flight controller fuses data from several specialized sensors. Science sensor can generate noisy data, the Flight Controller takes data from multiple sensors and blends them to find "optimal estimate". Most common sensors used for position estimation are:

\colger

 TBD

\coleng

\begin{itemize}
  \item \textsl{IMU (Inertial Measurement Unit)} - Tracks high-speed changes in acceleration and rotation in six directions to maintain basic flight attitude and stability.
  \item \textsl{GNSS (GPS)} - Establishes absolute global coordinates (latitude and longitude) for long-range navigation and geofencing.
  \item \textsl{Optical Flow} - Measures horizontal displacement by analyzing ground texture movement, enabling Position Hold in GNSS-denied (indoor) environments.
  \item \textsl{Barometer} - Measures air pressure to estimate relative altitude. While less precise than LiDAR at low altitudes, it is essential for maintaining height at scales beyond the rangefinder's limit.
  \item \textsl{LiDAR (Rangefinder)} - Provides high-precision at low-altitude distance measurements via laser pulses, critical for terrain following and Precision Landing.
\end{itemize}

\colger

\begin{itemize}
  \item \textsl{IMU (Inertial Measurement Unit)} - TBD
  \item \textsl{GNSS (GPS)} - TBD
  \item \textsl{Optical Flow} - TBD
  \item \textsl{Barometer} - TBD
  \item \textsl{LiDAR (Rangefinder)} - TBD
\end{itemize}

\coleng

With the accurate estimation of the drone’s state, the flight controller can move beyond manual flight into autonomous modes. In these modes, the system continuously monitors the drone’s position and automatically corrects any drift using sensor feedback:

\colger

TBD

\coleng

\begin{itemize}
  \item \textsl{Altitude Hold (Z-axis)} - The drone maintains a fixed height by using a LiDAR sensor at low altitude and a barometer at a higher altitude. This allows the operator to release the throttle while the drone keeps hovering at the same height. As this mode can be used with a barometer only and most modern flight controller comes with barometer preinstalled, Betflight, iNav and Arudpilot can activate this mode without requireing extra sensor.
  \item \textsl{Position Hold (X, Y axes)} - The drone keeps its horizontal position by using optical flow or GNSS (GPS) to detect movement. If wind or another force pushes the drone, the flight controller tilts the drone slightly to move it back to the same spot.
  \item \textsl{Return to Home (RTH / RTL)} - This is one of the most common safety features in modern drones, where the drone automatically flies back to its recorded home location using GNSS coordinates if the radio connection is lost or the battery becomes low.
\end{itemize}

\colger

\begin{itemize}
  \item \textsl{Altitude Hold (Z-axis)} - TBD
  \item \textsl{Position Hold (X, Y axes)} - TBD
  \item \textsl{Return to Home (RTH / RTL)} - TBD
\end{itemize}

\colende

\renewcommand{\deutschertitel}{GPS-basierte Positionsbestimmung}
\renewcommand{\englischertitel}{GPS-Based Position Estimation}
\makroabschnitt
\label{GpsBasedPosition}

The main method used for high-level autonomous navigation in drones is Global Navigation Satellite Systems (GNSS). GNSS is a general term that includes several satellite systems such as Global Positioning System (GPS) from the United States, GLONASS from Russia, and Galileo from the European Union. These satellite systems provide the unmanned aerial vehicle (UAV) with a global reference for positioning. Using these signals, the flight controller can estimate the drone’s absolute position in terms of latitude, longitude, and altitude.

\colger

TBD

\coleng

Modern GNSS receivers (see Section~\ref{AbschnittReceiver}) used in drones and flight controllers such as ArduPilot, Betflight, iNav and px4 can use signals from multiple satellite navigation systems at the same time (i.e., GPS + Galileo or GPS + GLONASS + Galileo). Using several systems together increases the number of visible satellites, which improves positioning accuracy and makes navigation more reliable, especially in challenging environments.

\colger

TBD

\colende

\begin{table}[htb!]
\centering
\caption{Firmware Configuration: From Hardware Setup to Mode Activation}
\label{TableFirmwareConfigForGPS}
\setlength{\tabcolsep}{4pt}
\renewcommand{\arraystretch}{1.4}
\scriptsize
\begin{tabularx}{\textwidth}{lXX}
\toprule
Firmware & Hardware Configuration (Enabling GPS) & Mode Activation (AltHold \& PosHold) \\
\midrule

Betaflight &
\begin{enumerate}
  \item Ports Tab: Set the specific UART to Sensor Input: GPS (see Figure~\ref{AbbildungBetaflightAppPorts}).
  \item GPS Tab(Configuration Tab in older betaflight app): Enable GPS and select the protocol (Most modern GPS module support UBLOX).
\end{enumerate}
&
\begin{enumerate} 
  \item Modes Tab: Assign a switch to ALTHOLD, POS HOLD and GPS RESCUE (see Figure~\ref{AbbildungBetaflightAppModes}).
  \item Failsafe Tab: Here GPS RESCUE mode can be set but other parameters of GPS RESCUE mode must be set carefully.
\end{enumerate} \\

\midrule

iNav &
\begin{enumerate} 
  \item Ports Tab: Enable GPS on the assigned UART.
  \item Configuration Tab: Enable GPS and Magnetometer.
  \item Perform a 3D compass calibration for navigation.
\end{enumerate}
&
\begin{enumerate} 
  \item Modes Tab: Assign switches to NAV ALTHOLD, NAV POSHOLD and NAV RTH.
\end{enumerate}  \\

\midrule

ArduPilot &
\begin{enumerate} 
  \item Setup Menu: The GPS is usually auto-detected on the GPS port and if not auto-detected set SERIALx\_PROTOCOL to 5 (here x is the serial number in the Flight controller board) (see Figure~\ref{MissionPlannerSerialPortConfig}) 
\end{enumerate}
&
\begin{enumerate} 
  \item Flight Modes: First assign a transmitter channel by using FLTMODE\_CH from full parameter list and  and from Setup tab → flight mode page, set AltHold, LOITER and RTL (see Figure~\ref{MissionPlannerFlightMode}).
\end{enumerate}
\\

\bottomrule
\end{tabularx}
\end{table}


\renewcommand{\deutschertitel}{Positionsbestimmung ohne GPS (Optischer Fluss und Entfernungsmesser)}
\renewcommand{\englischertitel}{GPS-Denied Position Estimation (Optical Flow and Lidar sensor)}
\makroabschnitt
\label{GpsDeniedPositionEstimation}

In environments such as indoor spaces where Global Navigation Satellite Systems (GNSS) signals are not available, drones can use optical flow and LiDAR sensors for state estimation. These sensors provide local measurements that help the flight controller determine the drone’s motion and altitude without using satellite signals. Flight control firmware such as Betaflight, iNav, and ArduPilot can use this sensor data to support automated flight modes such as Position Hold and Altitude Hold.

\colger

TBD

\colende

\renewcommand{\deutschertitel}{Optischer Fluss und Entfernungsmesser (Betaflight)}
\renewcommand{\englischertitel}{Opticalflow and LiDAR sensor Configuration (Betaflight)}
\makrounterabschnitt

Betaflight (as of version 4.6) is currently integrating support, primarily focusing on sensor data acquisition as a foundation for future autonomous features. To configure these sensors, the appropriate UART port must first be assigned to the corresponding peripheral in the Ports tab (see Figure~\ref{AbbildungBetaflightAppPorts}). Subsequently, the correct rangefinder and optical flow sensor types must be selected in the Configuration tab to enable their operation (see Figure~\ref{ImageLevelBetaflightAppLiderConfig}).

\colger

TBD

\colende

\begin{figure}[htb!]
  \centering
    \includegraphics[width=\linewidth]{Betaflight_2025.12.2_sensor_configuration.png}
  \caption{Configuration Page of Betaflight}
  \label{ImageLevelBetaflightAppLiderConfig}
\end{figure}
\FloatBarrier

\renewcommand{\deutschertitel}{Optischer Fluss und Entfernungsmesser (ArduPilot)}
\renewcommand{\englischertitel}{Opticalflow and LiDAR sensor Configuration (ArduPilot)}
\makrounterabschnitt

For reliable use of optical flow and rangefinder sensors, the drone must be capable of stable flight with minimal vibration. Excessive vibrations can significantly affect the accuracy of optical flow measurements and rangefinder readings, leading to unreliable sensor data and very bad flight (See how vibration can be reduced in ArduPilot~\ref{ArdupilotVibrationReduction}).

\colger

TBD

\coleng

Although ArduPilot supports installing optical flow and rangefinder sensors in different orientations, this section focuses only on the downward-facing configuration, since the installation and configuration procedures for other orientations are similar. It is recommended to mount the sensor close to the center of the drone and ensure that its forward direction is correctly aligned (see Figure~\ref{DroenWithMicoAirMFT01}). If this placement is not possible, the sensor position can be compensated through configuration parameters

\colende

% \begin{figure}[htbp]
%   \centering
%   \begin{minipage}[t]{0.48\textwidth}
%     \centering
%     \includegraphics[width=\linewidth]{MicoAir_Optical_Flow_Lidar_Sensor_MTF_01_in_flying_drone.jpg}
%     \vspace{0pt} % sorgt für Top-Ausrichtung
%   \end{minipage}\hfill
%   \begin{minipage}[t]{0.48\textwidth}
%     \centering
%     \includegraphics[width=\linewidth]{MicoAir_Optical_Flow_Lidar_Sensor_MTF_01_in_flying_drone.jpeg}
%     \vspace{0pt}
%   \end{minipage}
%   \caption{Drone installed with MicoAir MTF-01 Optical Flow and Lidar Sensor}
%   \label{DroenFlyingWithMicoAirMFT01}
% \end{figure}
% \FloatBarrier


\begin{figure}[htbp]
  \centering
  \begin{minipage}[t]{0.48\textwidth}
    \centering
    \includegraphics[width=\linewidth]{MicoAir_Optical_Flow_Lidar_Sensor_MTF_01_in_flying_drone.jpeg}
    \vspace{0pt} % sorgt für Top-Ausrichtung
  \end{minipage}\hfill
  \begin{minipage}[t]{0.48\textwidth}
    \centering
    \includegraphics[width=\linewidth]{MicoAir_Optical_Flow_Lidar_Sensor_MTF_01_in_drone.jpeg}
    \vspace{0pt}
  \end{minipage}
  \caption{Drone installed with MicoAir MTF-01 Optical Flow and Lidar Sensor}
  \label{DroenWithMicoAirMFT01}
\end{figure}
\FloatBarrier

